\documentclass[11pt]{article}

\usepackage{amsmath}
\usepackage{amssymb}
\usepackage{graphicx}
\usepackage{enumitem}
\usepackage[margin=1in]{geometry}
\usepackage[colorlinks=true, urlcolor=blue]{hyperref}

% Put the plots saved by the notebooks (the plots/ folder) next to this file,
% or place this file inside your HW4 repo folder.
\graphicspath{{plots/}{../plots/}{./}}

\newcommand{\answer}{\par\noindent\textbf{Answer:} }

\title{AA274A Homework 4}
\author{Name: \\ SUNet ID: }
\date{}

\begin{document}

\begin{document}
\maketitle

\noindent\textbf{Collaborators:} None

\noindent\textbf{Use of generative AI:} 

% Keep each problem on its own page and every part in the order below.

%%%%%%%%%%%%%%%%%%%%%%%%%%%%%%%%%%%%%%%%%%%%%%%%
\newpage
\section*{Problem 1: Convolutions and Image Features}

\subsection*{(i)}
Graded by the autograder; nothing to include.

\subsection*{(ii)}
% \begin{center}
%     \includegraphics[width=0.6\linewidth]{p1_matches.png}
% \end{center}

\answer Corners in first view: \underline{\hspace{1.5cm}}; corners in second view: \underline{\hspace{1.5cm}}; correct matches: \underline{\hspace{1.5cm}}

\subsection*{(iii)}

\subsection*{(iv)}
% \begin{center}
%     \includegraphics[width=0.6\linewidth]{p1_rotation.png}
% \end{center}

\answer
\begin{center}
\begin{tabular}{c|c|c}
Rotation & Repeatability & Correct matches \\ \hline
$0^\circ$  & \underline{\hspace{1.5cm}} & \underline{\hspace{1.5cm}} \\
$30^\circ$ & \underline{\hspace{1.5cm}} & \underline{\hspace{1.5cm}} \\
$90^\circ$ & \underline{\hspace{1.5cm}} & \underline{\hspace{1.5cm}}
\end{tabular}
\end{center}

\subsection*{(v)}
\noindent\textbf{(a)}

\answer Maximum difference away from the border: \underline{\hspace{1.5cm}}

\medskip
\noindent\textbf{(b)}

\answer Direct $13\times 13$ Gaussian: \underline{\hspace{1.5cm}} multiplications per pixel; separable: \underline{\hspace{1.5cm}} multiplications per pixel

\medskip
\noindent\textbf{(c)}

\answer Weights in one $3\times 3$ mask: \underline{\hspace{1.5cm}}; receptive field after three $3\times 3$ convolutions: \underline{\hspace{1.5cm}}; layers needed for a $15\times 15$ receptive field: \underline{\hspace{1.5cm}}

%%%%%%%%%%%%%%%%%%%%%%%%%%%%%%%%%%%%%%%%%%%%%%%%
\newpage
\section*{Problem 2: Bundle Adjustment SLAM with GTSAM}

\subsection*{(i)}
% Insert your labeled factor graph here (drawn by hand or with a drawing tool).
% \begin{center}
%     \includegraphics[width=0.75\linewidth]{p2_factor_graph.png}
% \end{center}

\answer Number of factors (including the prior): \underline{\hspace{1.5cm}}

\subsection*{(ii)}
Graded by the autograder; nothing to include.

\subsection*{(iii)}
% \begin{center}
%     \includegraphics[width=0.8\linewidth]{p2_slam_5.png}
% \end{center}

\answer Total factors: \underline{\hspace{1.5cm}}; prior factors: \underline{\hspace{1.5cm}}; odometry factors: \underline{\hspace{1.5cm}}; range-bearing factors: \underline{\hspace{1.5cm}}; loop-closure factors: \underline{\hspace{1.5cm}}

\answer Maximum position error (dead reckoning): \underline{\hspace{1.5cm}} m; maximum position error (optimized): \underline{\hspace{1.5cm}} m

\subsection*{(iv)}

\subsection*{(v)}
% \begin{center}
%     \includegraphics[width=0.8\linewidth]{p2_ranges.png}
% \end{center}

\answer
\begin{center}
\begin{tabular}{c|c|c}
Sensor range & Number of factors & Maximum position error (m) \\ \hline
$3$ m  & \underline{\hspace{1.5cm}} & \underline{\hspace{1.5cm}} \\
$5$ m  & \underline{\hspace{1.5cm}} & \underline{\hspace{1.5cm}} \\
$10$ m & \underline{\hspace{1.5cm}} & \underline{\hspace{1.5cm}}
\end{tabular}
\end{center}

% Report the final-pose uncertainty (standard deviation) with and without
% loop closure for each sensor range; use the same units as the notebook.
\answer
\begin{center}
\begin{tabular}{c|c|c}
Sensor range & Final-pose std. with loop closure & Without loop closure \\ \hline
$3$ m  & \underline{\hspace{1.5cm}} & \underline{\hspace{1.5cm}} \\
$5$ m  & \underline{\hspace{1.5cm}} & \underline{\hspace{1.5cm}} \\
$10$ m & \underline{\hspace{1.5cm}} & \underline{\hspace{1.5cm}}
\end{tabular}
\end{center}

\end{document}

\end{document}
